\documentclass[10pt,a4paper]{amsart}
\usepackage{amsmath,amssymb,amsthm,graphicx}
% Retained minimal preamble for the Stacks Project chapter more-algebra.tex (More on Algebra).
% The chapter file pinned-original-source.tex carries only \input{preamble} on its first line, so no
% preamble text lives inside the chapter itself.  The Stacks Project uses one single shared preamble for
% every chapter; the snapshot stages it identically as sources/stacks-chow/preamble.tex and
% sources/stacks-etale-cohomology/preamble.tex.  The block below is copied mechanically, line for line,
% from sources/stacks-chow/preamble.tex lines 14-18 (the Xy-pic setup) and 154-217 (the native theorem
% declarations together with the mathematical macros).  Nothing was added, removed or redefined; the
% wrapper class and the amsmath/amssymb/amsthm/graphicx base packages are supplied by the assembler.
% Chapter: more-algebra.tex -- More on Algebra.
% For commutative diagrams we use Xy-pic
\usepackage[all]{xy}

% We use 2cell for 2-commutative diagrams.
\xyoption{2cell}

% Theorem environments.
%
\theoremstyle{plain}
\newtheorem{theorem}[subsection]{Theorem}
\newtheorem{proposition}[subsection]{Proposition}
\newtheorem{lemma}[subsection]{Lemma}

\theoremstyle{definition}
\newtheorem{definition}[subsection]{Definition}
\newtheorem{example}[subsection]{Example}
\newtheorem{exercise}[subsection]{Exercise}
\newtheorem{situation}[subsection]{Situation}

\theoremstyle{remark}
\newtheorem{remark}[subsection]{Remark}
\newtheorem{remarks}[subsection]{Remarks}

\numberwithin{equation}{subsection}

% Macros
%
\def\lim{\mathop{\mathrm{lim}}\nolimits}
\def\colim{\mathop{\mathrm{colim}}\nolimits}
\def\Spec{\mathop{\mathrm{Spec}}}
\def\Hom{\mathop{\mathrm{Hom}}\nolimits}
\def\Ext{\mathop{\mathrm{Ext}}\nolimits}
\def\SheafHom{\mathop{\mathcal{H}\!\mathit{om}}\nolimits}
\def\SheafExt{\mathop{\mathcal{E}\!\mathit{xt}}\nolimits}
\def\Sch{\mathit{Sch}}
\def\Mor{\mathop{\mathrm{Mor}}\nolimits}
\def\Ob{\mathop{\mathrm{Ob}}\nolimits}
\def\Sh{\mathop{\mathit{Sh}}\nolimits}
\def\NL{\mathop{N\!L}\nolimits}
\def\CH{\mathop{\mathrm{CH}}\nolimits}
\def\proetale{{pro\text{-}\acute{e}tale}}
\def\etale{{\acute{e}tale}}
\def\QCoh{\mathit{QCoh}}
\def\Ker{\mathop{\mathrm{Ker}}}
\def\Im{\mathop{\mathrm{Im}}}
\def\Coker{\mathop{\mathrm{Coker}}}
\def\Coim{\mathop{\mathrm{Coim}}}

% Boxtimes
%
\DeclareMathSymbol{\boxtimes}{\mathbin}{AMSa}{"02}

%
% Macros for moduli stacks/spaces
%
\def\QCohstack{\mathcal{QC}\!\mathit{oh}}
\def\Cohstack{\mathcal{C}\!\mathit{oh}}
\def\Spacesstack{\mathcal{S}\!\mathit{paces}}
\def\Quotfunctor{\mathrm{Quot}}
\def\Hilbfunctor{\mathrm{Hilb}}
\def\Curvesstack{\mathcal{C}\!\mathit{urves}}
\def\Polarizedstack{\mathcal{P}\!\mathit{olarized}}
\def\Complexesstack{\mathcal{C}\!\mathit{omplexes}}
% \Pic is the operator that assigns to X its picard group, usage \Pic(X)
% \Picardstack_{X/B} denotes the Picard stack of X over B
% \Picardfunctor_{X/B} denotes the Picard functor of X over B
\def\Pic{\mathop{\mathrm{Pic}}\nolimits}
\def\Picardstack{\mathcal{P}\!\mathit{ic}}
\def\Picardfunctor{\mathrm{Pic}}
\def\Deformationcategory{\mathcal{D}\!\mathit{ef}}
\usepackage{cleveref}
\crefname{originallemma}{original lemma}{original lemmas}
\Crefname{originallemma}{Original lemma}{Original lemmas}
\makeatletter
\newlabel{lemma-prepare}{{09F2}{source}{}{sourcepointer-1}{}}
\newlabel{lemma-prepare@cref}{{[originallemma][0][]09F2}{[1][0][]source}{}{}{}}
\makeatother
\crefname{originallemma}{original lemma}{original lemmas}
\Crefname{originallemma}{Original lemma}{Original lemmas}
\makeatletter
\newlabel{lemma-characteristic-p-case}{{09F1}{source}{}{sourcepointer-2}{}}
\newlabel{lemma-characteristic-p-case@cref}{{[originallemma][0][]09F1}{[1][0][]source}{}{}{}}
\makeatother
\crefname{originallemma}{original lemma}{original lemmas}
\Crefname{originallemma}{Original lemma}{Original lemmas}
\makeatletter
\newlabel{lemma-cohen}{{09EZ}{source}{}{sourcepointer-3}{}}
\newlabel{lemma-cohen@cref}{{[originallemma][0][]09EZ}{[1][0][]source}{}{}{}}
\makeatother
\crefname{originallemma}{original lemma}{original lemmas}
\Crefname{originallemma}{Original lemma}{Original lemmas}
\makeatletter
\newlabel{lemma-extension-defined-by-nice-polynial}{{09F3}{source}{}{sourcepointer-4}{}}
\newlabel{lemma-extension-defined-by-nice-polynial@cref}{{[originallemma][0][]09F3}{[1][0][]source}{}{}{}}
\makeatother
\crefname{originallemma}{original lemma}{original lemmas}
\Crefname{originallemma}{Original lemma}{Original lemmas}
\makeatletter
\newlabel{lemma-weakly-unramified-goes-up-along-totally-ramified}{{09ER}{source}{}{sourcepointer-5}{}}
\newlabel{lemma-weakly-unramified-goes-up-along-totally-ramified@cref}{{[originallemma][0][]09ER}{[1][0][]source}{}{}{}}
\makeatother
\title{Lowering the level of a degree p extension of discretely valued fields by a totally ramified separable base change}
\author{The Stacks Project Authors}
\date{stacks-project:more-algebra.tex:local-snapshot-20261012; GNU FDL 1.2 or later}
\begin{document}
\maketitle
\setcounter{subsection}{1}
\noindent\textit{Editorial note.} All delivered mathematical text is the author's own, copied from the pinned Stacks Project chapter file more-algebra.tex; no statement or proof step was rewritten, reordered or removed.\par\medskip
\noindent\textit{Editorial note.} The chapter file pinned-original-source.tex carries only an input of the shared preamble on its first line, so no preamble text lives inside the chapter. The retained file sources/stacks-more-algebra/retained-preamble.tex is a mechanical, line-for-line copy of the project's single shared preamble as staged in the snapshot, namely sources/stacks-chow/preamble.tex lines 14-18 (Xy-pic setup) and 154-217 (theorem declarations and mathematical macros). Nothing was added, removed or redefined, no package was added, and the excerpt serves only as preamble, never as a span coordinate.\par\medskip
\noindent\textit{Editorial note.} One context span is delivered ahead of the statement: source lines 35142-35162, the author's own paragraph defining the level of a degree p extension of finite level. That paragraph is the only place in the chapter where the invariant appearing in the statement and in the conclusion is defined, and it is therefore shipped verbatim rather than paraphrased; it contains two forward pointers, to lemma-prepare and to lemma-extension-defined-by-nice-polynial, which are covered by the external reference stubs below. The same-chapter text that precedes it at source line 35140 is the end of another proof and is kept out of every counted span.\par\medskip
\noindent\textit{Editorial note.} The source URL is the source-level Stacks Project URL, as required when more than one card is produced from one source; the per-claim tags are carried by claim\_key and by the external references. All five external references are same-chapter pointers into the section Eliminating ramification and its immediate surroundings, recorded as non-load-bearing navigation stubs: lemma-prepare (09F2), lemma-characteristic-p-case (09F1), lemma-cohen (09EZ), lemma-extension-defined-by-nice-polynial (09F3) and lemma-weakly-unramified-goes-up-along-totally-ramified (09ER). The two equations labelled equation-first-congruence and equation-second-congruence are defined inside the counted proof span and resolve locally.\par\medskip
\noindent\textit{Editorial note.} Two obstacles in the source text are reproduced as they stand rather than repaired, because the author is responsible for the mathematics: at source line 35222 the exponent in the error term is printed with a stray closing parenthesis, and at source line 35230 the displayed definition of the set J uses the index i instead of j. Neither affects compilation and both belong to the original.\par\medskip
\noindent\textit{Allowed standard prerequisites (editorial).}\begin{itemize}
\item The definition of the level of a degree p extension of finite level, as fixed in the section immediately preceding the statement. the author's own paragraph at source lines 35142-35162, shipped as the counted context span of this card. Fixes the invariant that the statement is about: for a degree p extension L/K presented by an element z with xi = P(z) in K and pi\textasciicircum{}n xi integral, the level is the smallest value of n/e\_1, and it lies in the interval from 0 to p.
\item The normal form of a degree p extension of finite level over a discretely valued field: the polynomial P and the element w prescribed by the preparation lemma. Algebra chapter of this file, Lemma lemma-prepare (canonical tag 09F2), invoked at context line 35145 and proof line 35184. Supplies the polynomial P and the element w whose order e\_1 converts the integrality exponent n into the level l = n/e\_1, and whose multiplicative law P(x + y + wxy) = P(x) + P(y) + w\textasciicircum{}p P(x)P(y) is used at every correction step of the proof.
\item Cohen's theorem producing compatible residue-field representatives modulo a power of a uniformizer. Algebra chapter of this file, Lemma lemma-cohen (canonical tag 09EZ), invoked at proof lines 35214-35218. Produces the set-theoretic section sigma of the residue map used to expand xi in the truncated quotient and, later, to correct z.
\item The dictionary between the polynomial P(y) with constant term -c and the ramification behaviour of the resulting degree p extension, in the three cases (A), (B), (C). Algebra chapter of this file, Lemma lemma-extension-defined-by-nice-polynial (canonical tag 09F3), invoked at proof lines 35243, 35256, 35261 and 35365. Converts each verification about the value P(z') into a statement about separability, total ramification, or the field-versus-product alternative for the base-changed extension.
\item Totally ramified base change preserves the property of being weakly unramified. Algebra chapter of this file, Lemma lemma-weakly-unramified-goes-up-along-totally-ramified (canonical tag 09ER), invoked at proof lines 35260 and 35323. Keeps the hypothesis A in B weakly unramified available after the base changes K\_1/K that the proof performs.
\end{itemize}\medskip
\ifdefined\hypertarget\hypertarget{sourcepointer-1}{}\fi
\noindent\textit{Source reference.} Original lemma 09F2: sources/stacks-more-algebra/pinned-original-source.tex, source lines 34977--34977. Same-chapter pointer fixing the normal form of a degree p extension of finite level, namely the polynomial P and the element w; it is invoked at context line 35145 and proof line 35184 to name those objects, and the whole computation built on them stays inside the counted spans.\par\medskip
\ifdefined\hypertarget\hypertarget{sourcepointer-2}{}\fi
\noindent\textit{Source reference.} Original lemma 09F1: sources/stacks-more-algebra/pinned-original-source.tex, source lines 34818--34818. Same-chapter pointer mentioned once at proof lines 35191-35192 as the analogous equicharacteristic-p statement; the author explicitly states that the details of the present argument are given below, and every step of that argument is inside the counted proof span.\par\medskip
\ifdefined\hypertarget\hypertarget{sourcepointer-3}{}\fi
\noindent\textit{Source reference.} Original lemma 09EZ: sources/stacks-more-algebra/pinned-original-source.tex, source lines 34747--34747. Same-chapter pointer used at proof lines 35214-35218 to obtain compatible residue-field representatives modulo a power of the uniformizer; it supplies that section only, and the expansion, the index sets and the induction built on the section are all inside the counted proof span.\par\medskip
\ifdefined\hypertarget\hypertarget{sourcepointer-4}{}\fi
\noindent\textit{Source reference.} Original lemma 09F3: sources/stacks-more-algebra/pinned-original-source.tex, source lines 35017--35017. Same-chapter dictionary between the value of the polynomial P at a generator and the ramification behaviour of the resulting degree p extension, invoked at proof lines 35243, 35256, 35261 and 35365 to translate verifications that the counted proof performs explicitly.\par\medskip
\ifdefined\hypertarget\hypertarget{sourcepointer-5}{}\fi
\noindent\textit{Source reference.} Original lemma 09ER: sources/stacks-more-algebra/pinned-original-source.tex, source lines 34490--34490. Same-chapter pointer used at proof lines 35260 and 35323 to keep the weak unramifiedness hypothesis available after the two base changes; each base change is constructed explicitly in the counted proof span.\par\medskip
\noindent
Let $A$ be a discrete valuation ring of mixed characteristic $(0, p)$
containing a primitive $p$th root of $1$. Let $w \in A$ and $P(t) \in A[t]$
be as in Lemma \ref{lemma-prepare}. Let $L$ be a degree $p$ extension of
$K$ of finite level. Choose $z \in L$ generating $L$ over $K$
with $\xi = P(z) \in K$. Choose a uniformizer $\pi$ for $A$ and write
$w = u \pi^{e_1}$ for some integer $e_1 = \text{ord}_A(w)$
and unit $u \in A$. Finally, pick $n \geq 0$ such that
$$
\pi^n \xi \in A
$$
The {\it level} of $L/K$ is the smallest value of the quantity $n/e_1$
taking over all $z$ generating $L/K$ with $\xi = P(z) \in K$.

\medskip\noindent
We make a couple of remarks. Since the extension is of finite level
we know that we can choose $z$ such that $n < pe_1$.
Thus the level is a rational number contained in $[0, p)$.
If the level is zero then $L/K$ is unramified with respect to $A$ by
Lemma \ref{lemma-extension-defined-by-nice-polynial}.
Our next goal is to lower the level.

\begin{lemma}
\label{lemma-lowering-the-level}
Let $A \subset B \subset C$ be extensions of discrete valuation rings
with fractions fields $K \subset L \subset M$. Assume
\begin{enumerate}
\item $A$ has mixed characteristic $(0, p)$,
\item $A \subset B$ weakly unramified,
\item $B$ contains a primitive $p$th root of $1$,
\item $M/L$ is a degree $p$ extension of finite level $l > 0$,
\item $\kappa_A = \bigcap_{n \geq 1} \kappa_B^{p^n}$.
\end{enumerate}
Then there exists a finite separable extension $K_1$ of $K$
totally ramified with respect to $A$
such that either $K_1$ is a weak solution for $A \to C$, or the extension
$M_1/L_1$ is a degree $p$ extension of finite level
$\leq \max(0, l - 1, 2l - p)$.
\end{lemma}

\begin{proof}
Let $\pi \in A$ be a uniformizer.
Let $w \in B$ and $P \in B[t]$ be as in Lemma \ref{lemma-prepare} (for $B$).
Set $e_1 = \text{ord}_B(w)$, so that $w$ and $\pi^{e_1}$ are associates in $B$.
Pick $z \in M$ generating $M$ over $L$ with $\xi = P(z) \in K$
and $n$ such that $\pi^n\xi \in B$ as in the definition of the level
of $M$ over $L$, i.e., $l = n/e_1$.

\medskip\noindent
The proof of this lemma is completely similar to the proof of
Lemma \ref{lemma-characteristic-p-case}.
To explain what is going on, observe that
\begin{equation}
\label{equation-first-congruence}
P(z) \equiv z^p - z \bmod \pi^{-n + e_1}B
\end{equation}
for any $z \in L$ such that $\pi^{-n} P(z) \in B$ (use that $z$ has valuation
at worst $-n/p$ and the shape of the polynomial $P$). Moreover, we have
\begin{equation}
\label{equation-second-congruence}
\xi_1 + \xi_2 + w^p \xi_1 \xi_2 \equiv \xi_1 + \xi_2 \bmod \pi^{-2n + pe_1}B 
\end{equation}
for $\xi_1, \xi_2 \in \pi^{-n}B$. Finally, observe that
$n - e_1 = (l - 1)/e_1$ and $-2n + pe_1 = -(2l - p)e_1$.
Write $m = n - e_1 \max(0, l - 1, 2l - p)$. The above shows that doing
calculations in $\pi^{-n}B / \pi^{-n + m}B$ the polynomial $P$ behaves exactly
as the polynomial $z^p -  z$. This explains why the lemma is true
but we also give the details below.

\medskip\noindent
Assumption (4) implies that $\kappa_A$ is perfect. Observe that
$m \leq e_1$ and hence $A/\pi^m$ is annihilated by $w$ and hence $p$.
Thus we may choose compatible ring maps
$\overline{\sigma} : \kappa_A \to A/\pi^mA$ and
$\overline{\sigma} : \kappa_B \to B/\pi^mB$ as in
Lemma \ref{lemma-cohen}. We lift the second of these to a
map of sets $\sigma : \kappa_B \to B$. Then we can write
$$
\xi =
\sum\nolimits_{i = n, \ldots, n - m + 1} \sigma(\lambda_i) \pi^{-i} +
\pi^{-n + m)} b
$$
for some $\lambda_i \in \kappa_B$ and $b \in B$. Let
$$
I = \{i \in \{n, \ldots, n - m + 1\} \mid \lambda_i \in \kappa_A\}
$$
and
$$
J = \{j \in \{n, \ldots, n - m + 1\} \mid \lambda_i \not \in \kappa_A\}
$$
We will argue by induction on the size of the finite set $J$.

\medskip\noindent
The case $J = \emptyset$. Here for all $i \in \{n, \ldots, n - m + 1\}$
we have $\sigma(\lambda_i) = a_i + \pi^{n - m}b_i$ for some $a_i \in A$
and $b_i \in B$ by our choice of $\overline{\sigma}$. Thus
$\xi = \pi^{-n} a + \pi^{-n + m} b$ for some $a \in A$ and $b \in B$.
If $p | n$, then we write $a = a_0^p + \pi a_1$ for some $a_0, a_1 \in A$
(as the residue field of $A$ is perfect). Set $z_1 = - \pi^{-n/p} a_0$.
Note that $P(z_1) \in \pi^{-n}B$ and that $z + z_1 + w z z_1$ is an
element generating $M$ over $L$ (note that $wz_1 \not = -1$ as
$n < pe_1$). Moreover, by Lemma \ref{lemma-prepare} we have
$$
P(z + z_1 + w z z_1) = P(z) + P(z_1) + w^p P(z) P(z_1) \in K
$$
and by equations (\ref{equation-first-congruence}) and
(\ref{equation-second-congruence}) we have
$$
P(z) + P(z_1) + w^p P(z) P(z_1)
\equiv
\xi + z_1^p - z_1 \bmod \pi^{-n + m}B
$$
for some $b' \in B$. This contradict the minimality of $n$! Thus $p$
does not divide $n$. Consider the degree $p$ extension $K_1$ of $K$ given
by $P(y) = -\pi^{-n}a$. By Lemma \ref{lemma-extension-defined-by-nice-polynial}
this extension is separable and totally ramified with respect to $A$.
Thus $L_1 = L \otimes_K K_1$
is a field and $A_1 \subset B_1$ is weakly unramified
(Lemma \ref{lemma-weakly-unramified-goes-up-along-totally-ramified}).
By Lemma \ref{lemma-extension-defined-by-nice-polynial}
the ring $M_1 = M \otimes_K K_1$ is either a product of $p$ copies
of $L_1$ (in which case we are done) or a field extension of $L_1$
of degree $p$. Moreover, in the second case, either $C_1$ is weakly unramified
over $B_1$ (in which case we are done) or $M_1/L_1$ is degree $p$,
Galois, totally ramified with respect to $B_1$.
In this last case the extension $M_1/L_1$
is generated by the element $z + y + wzy$ and we see that
$P(z + y + wzy) \in L_1$ and
\begin{align*}
P(z + y + wzy)
& = P(z) + P(y) + w^p P(z) P(y) \\
& \equiv
\xi - \pi^{-n}a \bmod \pi^{-n + m}B_1 \\
& \equiv
0 \bmod \pi^{-n + m}B_1
\end{align*}
in exactly the same manner as above. By our choice of $m$ this
means exactly that $M_1/L_1$ has level at most $\max(0, l - 1, 2l - p)$.
From now on we assume that $J \not = \emptyset$.

\medskip\noindent
Suppose that $j', j \in J$ such that $j' = p^r j$ for some $r > 0$.
Then we set
$$
z_1 = - \sigma(\lambda_j) \pi^{-j} - \sigma(\lambda_j^p) \pi^{-pj} -
\ldots - \sigma(\lambda_j^{p^{r - 1}}) \pi^{-p^{r - 1}j}
$$
and we change $z$ into $z' = z + z_1 + wzz_1$. Observe that $z' \in M$
generates $M$ over $L$ and that we have
$\xi' = P(z') = P(z) + P(z_1) + wP(z)P(z_1) \in L$ with
$$
\xi' \equiv
\xi - \sigma(\lambda_j) \pi^{-j} + \sigma(\lambda_j^{p^r}) \pi^{-j'}
\bmod \pi^{-n + m}B
$$
by using equations (\ref{equation-first-congruence}) and
(\ref{equation-second-congruence}) as above. Writing
$$
\xi' = \sum\nolimits_{i = n, \ldots, n - m + 1} \sigma(\lambda'_i) \pi^{-i}
+ \pi^{-n + m}b'
$$
as before we find that
$\lambda'_i = \lambda_i$ for $i \not = j, j'$ and $\lambda'_j = 0$.
Thus the set $J$ has gotten smaller.
By induction on the size of $J$ we may assume there is no pair
$j, j'$ of $J$ such that $j'/j$ is a power of $p$.
(Please observe that in this procedure we may get thrown back into the case
that $J = \emptyset$ we treated above.)

\medskip\noindent
For $j \in J$ write $\lambda_j = \mu_j^{p^{r_j}}$ for some $r_j \geq 0$ and
$\mu_j \in \kappa_B$ which is not a $p$th power. This is possible by our
assumption (4). Let $j \in J$ be the unique index such that $j p^{-r_j}$
is maximal. (The index is unique by the result of the preceding paragraph.)
Choose $r > \max(r_j + 1)$ and such that $j p^{r - r_j} > n$ for $j \in J$.
Let $K_1/K$ be the extension of degree $p^r$, totally ramified
with respect to $A$, defined by $(\pi')^{p^r} = \pi$.
Observe that $\pi'$ is the uniformizer of the
corresponding discrete valuation ring $A_1 \subset K_1$.
Observe that $L_1 = L \otimes_K K_1$ is a field and $L_1/L$
is totally ramified with respect to $B$
(Lemma \ref{lemma-weakly-unramified-goes-up-along-totally-ramified}).
Computing in the integral closure $B_1$ we get
$$
\xi = \sum\nolimits_{i \in I} \sigma(\lambda_i) (\pi')^{-i p^r} +
\sum\nolimits_{j \in J} \sigma(\mu_j)^{p^{r_j}} (\pi')^{-j p^r} +
\pi^{-n + m} b_1
$$
for some $b_1 \in B_1$. Note that $\sigma(\lambda_i)$ for $i \in I$
is a $q$th power modulo $\pi^m$, i.e., modulo $(\pi')^{m p^r}$.
Hence we can rewrite the above as
$$
\xi = \sum\nolimits_{i \in I} x_i^{p^r} (\pi')^{-i p^r} +
\sum\nolimits_{j \in J} \sigma(\mu_j)^{p^{r_j}} (\pi')^{- j p^r}
+ \pi^{-n + m}b_1
$$
Similar to our choice in the previous paragraph we set
\begin{align*}
z_1 & - \sum\nolimits_{i \in I}
\left(x_i (\pi')^{-i} + \ldots + x_i^{p^{r - 1}} (\pi')^{-i p^{r - 1}}\right)
\\
& - \sum\nolimits_{j \in J}
\left(
\sigma(\mu_j) (\pi')^{- j p^{r - r_j}}
+ \ldots +
\sigma(\mu_j)^{p^{r_j - 1}} (\pi')^{- j p^{r - 1}}
\right)
\end{align*}
and we change our choice of $z$ into $z' = z + z_1 + wzz_1$.
Then $z'$ generates $M_1$ over $L_1$ and
$\xi' = P(z') = P(z) + P(z_1) + w^p P(z) P(z_1) \in L_1$
and a calculation shows that
$$
\xi' \equiv
\sum\nolimits_{i \in I} x_i (\pi')^{-i} +
\sum\nolimits_{j \in J} \sigma(\mu_j) (\pi')^{- j p^{r - r_j}} +
(\pi')^{(-n + m)p^r}b'_1
$$
for some $b'_1 \in B_1$. There is a unique $j$ such that $j p^{r - r_j}$
is maximal and $j p^{r - r_j}$ is bigger than $i \in I$. If
$j p^{r - r_j} \leq (n - m)p^r$ then the level of the extension $M_1/L_1$
is less than $\max(0, l - 1, 2l - p)$. If not, then, as $p$ divides
$j p^{r - r_j}$, we see that $M_1 / L_1$ falls into case (C) of
Lemma \ref{lemma-extension-defined-by-nice-polynial}.
This finishes the proof.
\end{proof}

\end{document}
